\documentclass[11pt,a4paper]{article}

\usepackage[T1]{fontenc}
\usepackage[utf8]{inputenc}
\usepackage[brazil]{babel}
\usepackage{lmodern}
\usepackage{microtype}
\usepackage[margin=2cm]{geometry}
\usepackage{hyperref}
\usepackage{titlesec}
\usepackage{xcolor}

\hypersetup{
  colorlinks=true,
  linkcolor=black,
  urlcolor=black,
  citecolor=black,
  pdfborder={0 0 0}
}

\definecolor{sectioncolor}{HTML}{111111}
\definecolor{accentcolor}{HTML}{1A1A1A}

\pagestyle{empty}
\setlength{\parindent}{0pt}
\setlength{\parskip}{5pt}
\linespread{1.05}
\selectfont

\titleformat{\section}
  {\bfseries\large\color{sectioncolor}}
  {}{0pt}{}[\vspace{-2pt}{\color{sectioncolor}\titlerule[0.8pt]}]
\titlespacing*{\section}{0pt}{12pt}{8pt}

\newcommand{\cvrole}[4]{%
  \vspace{4pt}%
  {\bfseries\color{accentcolor}#1}\hfill{\itshape #2}\\[-2pt]%
  {\itshape #3}\hfill #4\\[6pt]%
}

\begin{document}

\begin{center}
  {\fontsize{20}{24}\selectfont\bfseries Kamille Hartmann Maccagnan}\\[8pt]
  {\large PMO \textbar{} Gestão de Projetos \textbar{} Comunicação com Stakeholders}\\[10pt]
  Campo Comprido, Curitiba-PR \enspace\textbullet\enspace (41) 9 9928-4424 \enspace\textbullet\enspace
  \href{mailto:kamillehartmannm@gmail.com}{kamillehartmannm@gmail.com}\\[3pt]
  \href{https://linkedin.com/in/kamille-hartmann-maccagnan/}{linkedin.com/in/kamille-hartmann-maccagnan/}
\end{center}

\vspace{6pt}

\section{Resumo Profissional}

Profissional com experiência em gestão de projetos (PMO), acompanhamento de cronogramas e entregas, e comunicação com stakeholders em ambientes corporativos e multinacionais. Atua na organização de demandas, elaboração de relatórios executivos, monitoramento de indicadores de performance e padronização de processos, com foco em alinhamento entre áreas de interesses distintos e garantia de qualidade nas entregas. Vivência na interface com fornecedores, lideranças e equipes operacionais, incluindo reuniões multilíngues e atuação em contextos com múltiplas prioridades simultâneas. Experiência em documentação operacional, treinamento de equipes e disseminação de conhecimento em rede de lojas, com perfil analítico, organizado e orientado ao cumprimento de prazos.

\section{Experiência Profissional}

\cvrole{Anjuss}{Analista de Suporte de TI}{07/2025 -- Atual}

Atuo na organização de demandas internas, padronização de processos e apoio operacional à rede de lojas, com criação e atualização de documentações, POPs do sistema para as lojas e POPs voltados à equipe de T.I., assegurando consistência na execução e disseminação de informações entre as áreas. Lidero o treinamento e a capacitação de novos funcionários, promovendo alinhamento entre equipes e garantindo a adoção padronizada de procedimentos. Desenvolvo relatórios e dashboards em Power BI para monitoramento de indicadores operacionais e de desempenho, acompanhando a performance dos processos e produzindo visibilidade gerencial para apoio à tomada de decisão. Identifico oportunidades de melhoria operacional e automação de rotinas com Power Automate, contribuindo para maior eficiência, controle de qualidade das entregas e consistência das informações.

\cvrole{HORSE do Brasil}{Estagiária de TI (IS/IT LATAM)}{07/2024 -- 07/2025}

Atuei no acompanhamento de portfólio de projetos corporativos em ambiente multinacional, com participação em frentes de PMO, governança de dados, suporte e interface com áreas operacionais. Gerenciei o acompanhamento de cronogramas, prazos e entregas de iniciativas estratégicas, elaborando relatórios executivos e dashboards para monitoramento de performance e apoio à gestão de lideranças. Atuei como ponte entre áreas de negócio, fornecedores e stakeholders internacionais, participando de reuniões conduzidas em até três idiomas simultaneamente e promovendo alinhamento entre partes com demandas distintas. Utilizei o ServiceNow para gestão e acompanhamento de demandas, organização de chamados e coordenação entre usuários e áreas técnicas, e apliquei o Power Automate em iniciativas de automação de processos. Participei de iniciativas de melhoria contínua, análise de riscos, documentação de processos e evolução de práticas de governança de projetos.

\cvrole{Restaurante Familiar}{Analista de Dados e Suporte Técnico}{01/2023 -- 07/2024}

Atuei na análise de vendas, performance operacional e resultados financeiros para suporte à gestão do negócio. Estruturei do zero a base de dados em Excel para viabilizar dashboards em Power BI, desenvolvendo indicadores de desempenho para acompanhamento de resultados e produção de relatórios gerenciais. Realizei a padronização de informações para maior confiabilidade dos dados e identifiquei oportunidades de melhoria em processos operacionais. Atuei no alinhamento entre áreas operacionais e administrativas, apoiando a tomada de decisão com base em indicadores de performance, eficiência e qualidade dos processos.

\section{Competências Técnicas}

\textbf{PMO e Gestão de Projetos:} Acompanhamento de portfólio de projetos, controle de cronogramas e prazos, gestão de entregas, governança de projetos, gestão de demandas, análise de riscos, melhoria contínua, Scrum, Kanban, Boas Práticas PMOBOK.

\textbf{Comunicação e Stakeholders:} Interface entre áreas de negócio, fornecedores e lideranças, alinhamento entre stakeholders, comunicação com lideranças, reuniões multilíngues, mapeamento de necessidades, comunicação consultiva, gestão de relacionamento, treinamento e capacitação, onboarding de usuários.

\textbf{Indicadores e Relatórios:} Monitoramento de indicadores de performance, KPIs, relatórios executivos e gerenciais, dashboards em Power BI, análise de dados, Excel, SQL, acompanhamento de resultados operacionais e financeiros.

\textbf{Processos e Documentação:} Mapeamento de processos, padronização de processos, documentação operacional, POPs, elaboração de manuais e materiais de apoio, automação de processos (Power Automate).

\textbf{Ferramentas:} Power BI, ServiceNow, Trello, Monday, Spotfire, Pacote Office, Google Sheets.

\textbf{Idiomas:} Inglês Avançado \enspace\textbullet\enspace Espanhol Básico \enspace\textbullet\enspace Coreano Básico

\section{Projetos e Conquistas}

1º lugar no Transformation Day Renault 2023 com solução orientada a inovação, dados e melhoria de processos. Participação no Hackathon Bosch 2023 com foco em resolução de problemas de negócio. Desenvolvimento de dashboards executivos para monitoramento de indicadores de desempenho e suporte à tomada de decisão. Estruturação de KPIs para análise de performance operacional e acompanhamento de resultados. Automação e padronização de fluxos de dados e relatórios com Power Automate integrado ao Power BI, reduzindo retrabalho e aumentando confiabilidade das informações.

\section{Formação Acadêmica}

\textbf{PUCPR} \hfill Curitiba, PR\\
Graduação em Gestão da Tecnologia da Informação \hfill Previsão de conclusão: 06/2026\\[6pt]

\textbf{Escola Conquer} \hfill Curitiba, PR\\
Pós-graduação em Liderança e Gestão em Tecnologia \hfill Conclusão: 2024\\[6pt]

\textbf{Universidade Tuiuti do Paraná} \hfill Curitiba, PR\\
Graduação em Medicina Veterinária \hfill 06/2019\\[6pt]

\textbf{Qualittas} \hfill Curitiba, PR\\
Pós-graduação em Clínica Médica e Cirúrgica de Animais Exóticos e Pets Não Convencionais \hfill 12/2022

\end{document}
